\documentclass[11pt]{article}
\input{preamble}
\input{regrender}
\setdoc{Data}{Source discovery, data handling and audit protocol}{v0.1}{Draft -- EDA in progress}{2026-11-27}

\begin{document}
\makedoctitle
\tableofcontents
\clearpage

% =====================================================================
\section{Purpose}
This document has two parts.
Part 1 tells where the data comes from and which other sources exist.
Part 2 tells how we handle the data and how we audit its quality.
All numbers from the exploration notebook are \tObserved{} and preliminary (\refid{G-20}).

\taglegend

% =====================================================================
\part*{Part 1 -- Source discovery}
\addcontentsline{toc}{part}{Part 1 -- Source discovery}

\section{Why we search for sources}
The SNCF file is one output of a larger measurement system \tSource{} (search results, \refid{S-01}).
Other sources can answer our gaps, add stations, or contradict our framing (\refid{D-015}).
We search for sources that can refute our assumptions, not only for sources that support them.

\section{Search procedure}
\begin{enumerate}
  \item Trace the lineage of the dataset: SNCF portal $\rightarrow$ data.gouv.fr mirror $\rightarrow$ producer (SNCF Gares \& Connexions, Île-de-France) $\rightarrow$ measurement partners (Airparif, Agence d'Essai Ferroviaire) $\rightarrow$ coordinator (Île-de-France Mobilités).
  \item Search these portals: \url{https://ressources.data.sncf.com}, \url{https://www.data.gouv.fr}, \url{https://data.iledefrance-mobilites.fr}, \url{https://data.ratp.fr}, \url{https://www.airparif.fr}, \url{https://www.anses.fr}, \url{https://meteo.data.gouv.fr}.
  \item Search by station name (Avenue Foch, Sevran Beaudottes, La Défense, Neuilly -- Porte Maillot).
  \item Search by topic: ``qualité de l'air gares souterraines'', ``enceintes ferroviaires souterraines'', ``EFS''.
  \item Search for disconfirming evidence: sensor reliability problems, low exposure, no time pattern.
  \item Record each candidate in \fpath{registers/sources.csv}, also when you reject it. Write the reason.
  \item Set the decision: USE (this week), LATER (after delivery), REJECT (with reason).
\end{enumerate}

\section{Status of the search on 2026-11-27}
We ran three web searches. We tried to read the pages, but the network did not resolve the French sites from our environment.
So every source below is \sToCheck. The facts in the register come from search result summaries only. They are not read from the pages.
Read each page and update the register before you use a fact from it.

\textbf{First findings from search results} (to check):
\begin{itemize}
  \item SNCF Gares \& Connexions (Île-de-France) monitors air quality in underground rail spaces, with Airparif and the Agence d'Essai Ferroviaire. Avenue Foch (RER C) is one of four permanent stations (\refid{S-01}, \refid{S-03}).
  \item The data.gouv.fr page shows the licence as ``not specified'' (\refid{S-02}, \refid{G-13}).
  \item Île-de-France Mobilités published a map of particle levels for about 44 underground stations in 2024, with a yearly update, as open data (\refid{S-05}).
  \item No legal limit exists for these spaces. ANSES recommendations are used as a reference (\refid{S-05}, \refid{S-09}).
  \item RATP publishes hourly air data for some of its own stations (\refid{S-08}).
  \item Airparif reviewed the RATP monitoring system in 2022 and asked for more accessible and more detailed open data (\refid{S-07}).
\end{itemize}

\section{Source register}
\subsection{Summary}
\srcsummary
\subsection{Entries}
\srcentries

% =====================================================================
\part*{Part 2 -- Data handling and audit}
\addcontentsline{toc}{part}{Part 2 -- Data handling and audit}

\section{Data handling rule}
\sWorking{} (\refid{D-014}).
\begin{enumerate}
  \item The folder \fpath{spark/data/} is the raw zone. Do not write, edit, filter or overwrite files there.
  \item Add a file to the raw zone only with a documented download step.
  \item For each raw file, write a manifest entry in \fpath{registers/manifest.csv}.
  \item If a new download has a different checksum, keep both versions. Use a new file name with the download date.
  \item Write all processed data to \fpath{spark/lake/}:
  \begin{itemize}
    \item \fpath{bronze/}: typed copy of the raw file in Parquet, with ingestion time, source ID and checksum columns.
    \item \fpath{silver/}: same rows, with quality flag columns. No value is deleted or changed.
    \item \fpath{gold/}: features and labels for the models.
  \end{itemize}
  \item Make each Silver and Gold table reproducible from Bronze with code in the repository.
\end{enumerate}

\section{Raw file manifest}
\manifestentries

% =====================================================================
\section{Data dictionary}
The table gives each column of \fpath{qualite-de-lair.csv} and the status of each property.

\small
\begin{longtable}{@{}P{0.13\linewidth} P{0.11\linewidth} P{0.36\linewidth} P{0.32\linewidth}@{}}
\toprule
\textbf{Column} & \textbf{Notebook name} & \textbf{Meaning} & \textbf{Status of the properties} \\
\midrule\endhead
\bottomrule\endfoot
Trigramme & code & Three-letter station code (AVF, BDE, DGA, NPT) & Values \tObserved \\
Date heure & ts & Time of the measurement, ISO 8601 with suffix Z & Format \tObserved. Start or end of hour \tUnknown{} (\refid{G-02}) \\
PM10 & PM10 & \pmten{} concentration & Unit \ugm{} \tAssumed{} (\refid{G-03}). Averaging \tUnknown{} (\refid{G-01}) \\
PM25 & PM25 & \pmtwofive{} concentration & Same as PM10 \\
Nom\_Gare & station & Station name & Values \tObserved \\
Id\_Gare & station\_id & Station identifier (UUID) & Link to SNCF station referential \tUnknown{} (\refid{G-19}) \\
\end{longtable}
\normalsize

Personal data: none \tObserved. The file holds only sensor values, times and station names.

% =====================================================================
\section{Quality checks}
Run each check on Bronze. Write the result in Silver as a flag. Write the count per station in the notebook and here.

\small
\begin{longtable}{@{}P{0.15\linewidth} P{0.3\linewidth} P{0.27\linewidth} P{0.2\linewidth}@{}}
\toprule
\textbf{Dimension} & \textbf{Check} & \textbf{Spark method} & \textbf{Flag or gap} \\
\midrule\endhead
\bottomrule\endfoot
Uniqueness & One row for each station and hour & \texttt{groupBy(code, ts).count()} & \texttt{dup\_flag} \\
Completeness & Missing hours and null values per station; length of null runs & Window with \texttt{lag}, gaps-and-islands & \texttt{gap\_run\_h}, \refid{G-07} \\
Validity & Negative values & \texttt{filter(col < 0)} & \texttt{neg\_flag} \\
Consistency & \pmtwofive{} greater than \pmten & column comparison & \texttt{pm25\_gt\_pm10}, \refid{G-09} \\
Validity & Value equal to 5.0 & equality test, run length & \texttt{at\_floor}, \refid{G-08} \\
Plausibility & Extreme values (limits to fix with \refid{G-10}) & quantiles per station & \texttt{extreme\_flag} \\
Plausibility & Same value for many hours & run length of equal values & \texttt{stuck\_flag}, \refid{G-11} \\
Timeliness & Delay between last timestamp and download & manifest and \texttt{max(ts)} & \refid{G-12} \\
Coverage & Valid hours per station and year & \texttt{groupBy(code, year)} & \refid{G-18} \\
\end{longtable}
\normalsize

% =====================================================================
\section{Preliminary observations}
\textbf{EDA in progress -- not validated.} The values come from the saved outputs of the exploration notebook (\fpath{spark/apps/.ipynb_checkpoints/}), not from a fresh run (\refid{G-20}).
Each interpretation has at least one alternative.

\subsection{Size and coverage} \tObserved{} (notebook sections 1--2)
\begin{itemize}
  \item 206\,904 rows, from 2018-01-01 00:00 to 2026-11-25 23:00 (UTC as written in the file).
  \item No duplicate station and hour. No missing hourly timestamp (the rows exist; values can be null).
  \item Rows per station: AVF 78\,024, BDE 78\,024 (from 2018), NPT 34\,200 (from 2023), DGA 16\,656 (from 2025).
\end{itemize}

\subsection{Distributions} \tObserved{} (notebook section 3)
\begin{itemize}
  \item \pmten: mean 160.6, median 145.3, maximum 13\,633.79, skewness 10.38.
  \item \pmtwofive: mean 64.9, median 55.4.
  \item Correlation \pmten{} and \pmtwofive: 0.927. Median ratio \pmtwofive/\pmten: 0.41.
\end{itemize}

\subsection{Quality} \tObserved{} (notebook section 4)
\begin{itemize}
  \item Null values: \pmten{} 15.04\,\%, \pmtwofive{} 15.85\,\%. Between 89\,\% and 94\,\% of nulls are in runs of 24 hours or more. Longest run: 3\,761 hours (BDE).
  \item Negative values: 19 rows for \pmten, 19 rows for \pmtwofive{} (31 rows with one or both).
  \item \pmtwofive{} greater than \pmten: 2\,793 rows.
  \item Value exactly 5.0: 1\,085 rows (\pmten), 4\,149 rows (\pmtwofive).
  \item \pmten{} greater than 1000: 62 rows.
  \item BDE: \pmten{} equal to 5.0 for 121 hours, from 2024-12-27 to 2025-01-01.
\end{itemize}
\textbf{Interpretations and alternatives.}
\begin{itemize}
  \item 5.0 is a sensor floor value \tInterp. Alternative: a default value written by the data system during a fault. Gap \refid{G-08}.
  \item The 121 hours at 5.0 are a stuck sensor \tInterp. Alternative: a real low period during the holiday break with reduced service. Gap \refid{G-11}.
  \item Values above 1000 are faults \tInterp. Alternative: real events (works, incidents, cleaning). Gap \refid{G-10}.
\end{itemize}

\subsection{Time patterns} \tObserved{} (notebook section 5)
\begin{itemize}
  \item Weekend median \pmten{} is lower than weekday median: AVF $-29$\,\%, DGA $-28$\,\%, NPT $-46$\,\%, BDE $-17$\,\%.
\end{itemize}
\textbf{Interpretation and alternatives.}
The notebook reads this as a rail-traffic signal \tInterp.
Alternatives: ventilation schedule; outdoor road traffic; outdoor PM levels; maintenance or cleaning schedule; station occupancy (people move dust).
Hypothesis \refid{H-08} tests this. Gap \refid{G-16} lists the data needed.

\subsection{Candidate flag rules} \sProposed{} (\refid{D-009})
The notebook proposes these rules. They become flags, not deletions. Each depends on a gap.
\begin{itemize}
  \item Negative value $\rightarrow$ \texttt{neg\_flag}.
  \item \pmtwofive{} greater than \pmten{} $\rightarrow$ \texttt{pm25\_gt\_pm10} on both values (\refid{G-09}).
  \item Value 5.0 $\rightarrow$ \texttt{at\_floor} (\refid{G-08}).
  \item \pmten{} above 1000 or \pmtwofive{} above 500 $\rightarrow$ \texttt{extreme\_flag} (\refid{G-10}). The limits are not justified yet.
\end{itemize}

\subsection{Known issues in the notebook}
\begin{itemize}
  \item A cell runs the heat map for BDE (Sevran Beaudottes), but its comment says ``Avenue Foch''.
  \item The summary says that the data has no coordinates. The file \fpath{spark/data/stations.csv} now has coordinates (origin: \refid{G-21}).
  \item The main notebook has no saved outputs. The numbers above come from the checkpoint copy (\refid{G-20}).
\end{itemize}

\end{document}
